\chapter{The Mouse}
\label{ch:mouse}

\index{mouse}

This is the most important chapter in this book.  Acme is a
mouse-driven editor.  The three mouse buttons are the primary way you
interact with acme---more so than the keyboard.  If you learn nothing
else, learn the three buttons and the chords.

Throughout this chapter, the buttons are:

\begin{description}
\item[\button{1}] Left button --- \textbf{Select}
\item[\button{2}] Middle button --- \textbf{Execute}
\item[\button{3}] Right button --- \textbf{Look / Open}
\end{description}

\section{Button 1: Select}
\label{sec:b1}

\index{mouse!Button 1}
\index{selection}

\button{1} is for selecting text.

\begin{description}[style=nextline]
\item[Click]
Places the cursor (an empty selection) at the click position.

\item[Drag]
Selects a range of text.  The selected text is highlighted.

\item[Double-click]
Selects the word under the cursor.  A ``word'' is a run of
alphanumeric characters and underscores, or a run of
non-whitespace characters, depending on context.

\item[Double-click on a bracket]
If you double-click on any of \texttt{( ) [ ] \{ \} < >}, acme
selects the text enclosed by the matching pair, including nested
brackets.  This works for parentheses, square brackets, curly
braces, and angle brackets.

\item[Double-click on a quote]
Double-clicking on a single quote, double quote, or backtick
selects the text enclosed by the matching quote character.
\end{description}

Clicking \button{1} anywhere deselects the previous selection.  The
selection is fundamental to acme: many commands operate on the selected
text, and the selected text is what gets cut, copied, or replaced.

\section{Button 2: Execute}
\label{sec:b2}

\index{mouse!Button 2}
\index{executing commands}

\button{2} executes text as a command.

\begin{description}[style=nextline]
\item[Click on a built-in command]
Middle-click on \cmd{Put} and acme saves the file.  Middle-click on
\cmd{Del} and acme closes the window.  Middle-click on \cmd{Undo}
and acme undoes the last change.

\item[Click on arbitrary text]
Middle-click on any text to run it as a shell command.  For example,
type \cmd{date} in a tag bar, middle-click it, and the current date
appears in the \file{+Errors} window.

\item[Click on a blank area]
If you click in an empty area or on whitespace, acme expands the
selection to the surrounding non-whitespace text and executes that.
This means you do not need to carefully select a command word---just
click anywhere on it.

\item[Click on a pipe command]
Text beginning with \texttt{<}, \texttt{>}, or \texttt{|} has special
meaning.  See Chapter~\ref{ch:shell} for details.
\end{description}

The command executes in the directory associated with the window.  For a
file \file{/home/user/src/main.c}, commands run in
\file{/home/user/src/}.

\section{Button 3: Look and Open}
\label{sec:b3}

\index{mouse!Button 3}
\index{Look}
\index{plumbing}

\button{3} is the ``smart open'' button.  It examines the text you
click on and does the most useful thing:

\begin{description}[style=nextline]
\item[Filename]
Right-click on \file{main.c} and acme opens it in a new window (or
jumps to the existing window if it is already open).

\item[Filename with address]
Right-click on \file{main.c:42} and acme opens \file{main.c} and
jumps to line~42.  The address can be a line number, a line:column
pair (\file{main.c:42.10}), or a regular expression
(\file{main.c:/pattern/}).

\item[Address in current file]
Right-click on \texttt{:42} or \texttt{:/pattern/} to jump to that
location in the current window.

\item[Include file]
Right-click on \texttt{<stdio.h>} and acme searches include
directories for \file{stdio.h} and opens it.

\item[URL]
Right-click on \texttt{https://example.com} and acme opens it in
your web browser (via \cmd{xdg-open}, \cmd{open}, or
\cmd{cygstart} depending on platform).

\item[Plain text]
Right-click on any other text and acme searches forward in the
current window body for that text---a literal search (not a regex).
Click again to find the next occurrence.
\end{description}

The decision-making process is handled by the \emph{plumber} (see
Chapter~\ref{ch:plumbing}).  The built-in plumbing rules handle
filenames, URLs, images, and PDFs.  You can add custom rules.

\section{Mouse Chords}
\label{sec:chords}

\index{mouse!chords}
\index{chords}

Mouse chords are combinations of buttons pressed in sequence.  They
provide quick access to cut, paste, and copy without reaching for menu
commands.

\subsection{Cut: \chord{1}{2}}
\index{Cut!chord}

\begin{enumerate}
\item Press and hold \button{1}, drag to select text.
\item While still holding \button{1}, click \button{2}.
\item Release both buttons.
\end{enumerate}

The selected text is deleted and placed in the snarf buffer (clipboard).
This is identical to executing the \cmd{Cut} command.

\subsection{Paste: \chord{1}{3}}
\index{Paste!chord}

\begin{enumerate}
\item Press and hold \button{1} to position the cursor or select text.
\item While still holding \button{1}, click \button{3}.
\item Release both buttons.
\end{enumerate}

The snarf buffer contents replace the selection (or are inserted at the
cursor position).  This is identical to executing the \cmd{Paste}
command.

\subsection{Snarf: \chord{1}{2} then \button{3}}
\index{Snarf!chord}

\begin{enumerate}
\item Press and hold \button{1}, drag to select text.
\item While holding \button{1}, click \button{2} (this cuts).
\item While still holding \button{1}, click \button{3} (this pastes
      back what was just cut).
\item Release all buttons.
\end{enumerate}

The net effect: the text remains unchanged, but it has been copied to
the snarf buffer.  This is the chord for ``copy without deleting.''

\subsection{The 2-1 Argument Chord}
\index{2-1 chord}

Commands can receive an argument via a two-button chord:

\begin{enumerate}
\item Press and hold \button{2} on a command word (e.g., \cmd{Look}).
\item While holding \button{2}, click \button{1} on the argument text.
\item Release both buttons.
\end{enumerate}

Acme passes the \button{1}-selected text as an argument to the
\button{2} command.  For example, holding \button{2} on \cmd{Look}
and clicking \button{1} on a word searches for that word.

\section{The Scrollbar}
\label{sec:scrollbar}

\index{scrollbar}

The scrollbar is the narrow strip along the left edge of each window
body.  It contains a thumb (small rectangle) showing the current
viewport position relative to the whole file.

\begin{description}[style=nextline]
\item[\button{1} in scrollbar]
Scrolls \textbf{up}.  The line at the click position moves to the
top of the window.

\item[\button{2} in scrollbar]
\textbf{Jump scroll}.  The display jumps so that the thumb is at the
click position---proportional scrolling.

\item[\button{3} in scrollbar]
Scrolls \textbf{down}.  The line at the top of the window moves to
the click position.
\end{description}

The \cmd{-r} command-line flag reverses the sense of \button{1} and
\button{3} in the scrollbar, for users who prefer the opposite
convention.

\section{Scroll Wheel}
\label{sec:scrollwheel}

\index{scroll wheel}

The scroll wheel scrolls the window body:

\begin{itemize}
\item Scroll up = scroll text up (view earlier content).
\item Scroll down = scroll text down (view later content).
\end{itemize}

The scroll amount is a fraction of the visible window height.

\section{One-Button and Two-Button Mice}
\label{sec:emulation}

\index{mouse!emulation}
\index{Ctrl-click}
\index{Alt-click}

On laptops, trackpads, and Mac mice, you may not have three distinct
buttons.  Acme provides modifier-click emulation:

\begin{longtable}{@{}ll@{}}
\toprule
Modifier Click & Equivalent \\
\midrule
\endhead
\key{Ctrl}-click & \button{2} (execute) \\
\key{Alt}-click (or \key{Option}-click on Mac) & \button{3} (look/open) \\
\bottomrule
\end{longtable}

These modifiers work with \button{1} (left click).  So on a trackpad:

\begin{itemize}
\item Plain click = \button{1} (select)
\item \key{Ctrl}+click = \button{2} (execute)
\item \key{Alt}+click = \button{3} (look/open)
\end{itemize}

\textbf{Recommendation:} For serious use, get a three-button mouse.
Acme was designed for one, and the experience is dramatically better
with real buttons.  Any inexpensive USB mouse with a clickable scroll
wheel provides three buttons.

\section{Summary}

\begin{longtable}{@{}llp{2.5in}@{}}
\toprule
Button & Name & Action \\
\midrule
\endhead
\button{1} & Select & Select text, position cursor, drag to move windows \\
\button{2} & Execute & Run built-in commands or shell commands \\
\button{3} & Look & Open files, search text, follow addresses \\
\chord{1}{2} & Cut & Delete selection to snarf buffer \\
\chord{1}{3} & Paste & Replace selection with snarf buffer \\
\chord{1}{2}+\button{3} & Snarf & Copy selection to snarf buffer \\
\bottomrule
\end{longtable}
